\documentclass[11pt]{article}
\usepackage[margin=1in]{geometry}
\usepackage[T1]{fontenc}
\usepackage[utf8]{inputenc}
\usepackage{lmodern}
\usepackage{hyperref}
\usepackage{parskip}

\title{Concise Summary of the Manuscript Revisions}
\author{}
\date{March 15, 2026}

\begin{document}
\maketitle

\section*{Overview}
The manuscript was revised to align its wording with regenerated evidence obtained from the cleaned,
auditable optimization pipeline. The revisions do not change the basic analytic construction of the
paper. Instead, they narrow several claims so that the text matches what the reproducible numerical
workflow now supports.

\section*{What Was Preserved}
The following points remain central to the paper:
\begin{itemize}
\item the work is a differentiable, physics-penalized optimization over a structured analytic ansatz,
\item the construction preserves non-negative seed energy density under the restricted uniform,
time-independent translations studied here,
\item the remaining obstruction is stress-dominated rather than density-dominated,
\item the double-shell geometry remains more constrained than the single-shell one,
\item the DEC remains the tightest residual condition in the regenerated examples.
\end{itemize}

\section*{What Was Softened}
The revised manuscript now avoids stronger formulations that were not supported by the regenerated
bundles or by direct evaluations of the manuscript target parameters. In particular, we no longer
state that:
\begin{itemize}
\item the NEC or WEC reach a high-success regime,
\item the sampled SEC slices are uniformly positive,
\item the visible main optimization run shows a large transient before plateauing,
\item the visible training dynamics are best described as a strong drop in physics loss.
\end{itemize}

\section*{Why the Changes Were Necessary}
After regenerating the single-shell and double-shell cases with the cleaned pipeline, and after
evaluating the manuscript parameter targets under the same diagnostics, the evidence consistently
showed the same qualitative hierarchy:
\begin{itemize}
\item $\rho$ remains non-negative by construction in the restricted setup considered,
\item the DEC is the hardest residual test,
\item the SEC is milder than the DEC in aggregate,
\item the NEC/WEC remain substantially constrained rather than nearly saturated.
\end{itemize}

The revised wording therefore reflects a deliberate effort to distinguish between the paper's
constructive analytic contribution and stronger energy-condition claims that the regenerated
numerical evidence does not justify.

\section*{Resulting Editorial Position}
The manuscript now presents a more limited but more defensible claim: selective warping of
positive-energy static seeds can be combined with reproducible constrained optimization so that the
remaining deficits are localized and stress-dominated, even though full pointwise energy-condition
satisfaction is not established.

\end{document}
